\ifx\XeTeXversion\undefined\pdfoutput=1\fi
\documentclass[11pt,a4paper]{article}

\usepackage[a4paper,margin=20mm]{geometry}
\usepackage{lmodern}
\usepackage{amsmath,amssymb,amsthm,mathtools}
\usepackage{booktabs}
\usepackage{xcolor}
\usepackage{microtype}
\usepackage{fancyhdr}
\usepackage{hyperref}
\usepackage{enumitem}
\usepackage{graphicx}
\usepackage[bottom]{footmisc}

\definecolor{primary}{RGB}{24,78,119}
\definecolor{accent}{RGB}{181,64,47}

\hypersetup{
    colorlinks=true,
    linkcolor=primary,
    citecolor=primary,
    urlcolor=primary
}

\setlength{\parskip}{6pt plus 2pt minus 1pt}
\setlength{\parindent}{0pt}

\newtheorem{theorem}{Theorem}[section]
\newtheorem{lemma}[theorem]{Lemma}
\newtheorem{proposition}[theorem]{Proposition}
\newtheorem{corollary}[theorem]{Corollary}

\theoremstyle{definition}
\newtheorem{definition}[theorem]{Definition}
\newtheorem{remark}[theorem]{Remark}

\newcommand{\Perm}{\mathrm{Perm}}
\newcommand{\rank}{\mathrm{rank}}
\newcommand{\F}{\mathbb{F}}
\newcommand{\R}{\mathbb{R}}
\newcommand{\Z}{\mathbb{Z}}
\newcommand{\NP}{\mathbf{NP}}
\newcommand{\sharpp}{\#\mathbf{P}}

\newcommand{\takeaway}[2]{%
  \vspace{6pt}\noindent\colorbox{primary!8}{%
    \parbox{\dimexpr\linewidth-2\fboxsep\relax}{%
      \vspace{3pt}\textbf{\color{primary}#1:}\hspace{0.5em}#2\vspace{3pt}%
    }%
  }\vspace{6pt}%
}

\pagestyle{fancy}
\fancyhf{}
\lhead{\small\color{gray}Sub-$2^n$ Sparse Matrix Permanent via Graph Reductions and Min-Frontier DP}
\rhead{\small\color{gray}\thepage}
\renewcommand{\headrulewidth}{0.4pt}

\title{%
    \textbf{Exact Sub-$2^n$ Matrix Permanent Evaluation via Iterated Series Reductions and Min-Frontier Path Decomposition}\\[0.4em]
    \large\textsf{Part II: Empirical Sub-$2^n$ Scaling on Bounded-Degree and Sparse Graphs}
}

\author{
    \textbf{Barath B}\thanks{Independent Researcher. Email: \texttt{barathjack77@gmail.com}.\\Correspondence: \url{https://github.com/3DBarath/matrix_permanent_circuit_discovery}.}
}
\date{November 2026}

\begin{document}

\maketitle

\begin{abstract}
Computing the permanent of an $n \times n$ matrix is canonical for Valiant's counting class $\sharpp$. While companion work on algebraic circuit discovery (Part~I) achieved $O(2^n)$ non-scalar multiplications (and $O(n \cdot 2^n)$ total arithmetic operations) for general dense matrices, all known dense algorithms remain constrained by the exponential base $2$. In this work, we show that for sparse, bounded-degree, and structured bipartite matrices, empirical scaling well below $2^n$ can be achieved without sacrificing exactness (without claiming a new worst-case asymptotic theorem across all graphs, and contrasting with proved worst-case bounds~\cite{servediowan2005,fominhoie2006,izumiwadayama2012,bjorklund2012,bjorklund2016}). We present an exact solver combining three components: (1) iterated zero-division graph reductions run to a fixed point, comprising forced degree-1 edge peeling, degree-2 series contractions that merge neighbor vertices algebraically in $O(\deg(c_1)+\deg(c_2))$ time ($O(1)$ when neighbor degrees are $O(1)$), and bipartite connected-component factorization with early Hall-condition detection; (2) a greedy min-frontier row-ordering heuristic targeting compact path decompositions of the underlying incidence graph, restricting dynamic programming state allocations to the active cut frontier; and (3) a repository extension providing hybrid dispatch to the Generation~9 Fused M\"obius--Zeta Butterfly OptTree engine for small dense components.

Our implementation is verified bit-exact against Glynn's baseline over 61-bit Mersenne primes $\F_p$ across all tested graph structures. In empirical fits on random 3-regular bipartite graphs (high-expansion stress tests, rather than a proved worst-case bound across all 3-regular graphs), the solver scales as $1.238^n$, evaluating $n = 28$ instances in $5{,}118$ DP operations ($52{,}446\times$ fewer than dense $2^n$) in $3.60$\,ms. On random 4-regular bipartite graphs, we observe an empirical scaling of $1.416^n$, computing $n = 24$ in $19.71$\,ms ($516.6\times$ fewer operations). Taken together, these results yield a practical, verified exact permanent solver that operates well below $2^n$ in practice without floating-point approximations.
\end{abstract}

\vspace{0.5em}
\noindent\textbf{Keywords:} Matrix Permanent, \#P-Complete, Sub-Exponential Graphs, Graph Contraction, Path Decomposition, Frontier Dynamic Programming.

\section{Introduction}

The matrix permanent of an $n \times n$ matrix $A$, defined by
\begin{equation}
\Perm(A) = \sum_{\sigma \in S_n} \prod_{i=1}^n A_{i, \sigma(i)},
\end{equation}
is a foundational counting problem in computational complexity~\cite{valiant1979}. In quantum optics, $\Perm(A)$ governs the transition amplitudes of indistinguishable photons in linear optical networks (Boson Sampling)~\cite{aaronson2011}, and in statistical mechanics it counts monomer-dimer covers and perfect matchings in bipartite graphs.

For arbitrary dense matrices, Ryser's classical inclusion-exclusion algorithm \cite{ryser1963} and Glynn's polarization formula \cite{glynn2010} established the baseline complexity of $(n-1)2^{n-1}$ non-scalar multiplications and $\Theta(n \cdot 2^n)$ total gates. In companion work (Part~I) \cite{part1_paper}, the author eliminated the polynomial factor $n$ from multiplicative complexity, achieving $O(2^n)$ non-scalar multiplications ($1.602 \cdot 2^{15}$ at $n = 16$) with $\Theta(n \cdot 2^n)$ total gates.

However, for dense matrices over general fields, the exponential base $2$ is bounded from below by the necessity of sampling the full Boolean hypercube $\{-1, 1\}^n$ to annihilate lower-degree cross-terms.

\paragraph{The Sparse Frontier and Related Work.}
In many applications---such as linear optical circuits with bounded optical depth, planar lattice models, and sparse matching problems---the incidence matrix $A$ is sparse. A matrix is termed $d$-regular if every row and column has exactly $d$ non-zero entries. On bounded-degree graphs, rigorous worst-case sub-$2^n$ bounds are already known: Servedio and Wan~\cite{servediowan2005} and bounded-average-degree DP~\cite{izumiwadayama2012} give $O^*(2^{(1 - 1/(3.55d))n})$ (roughly $1.87^n$ at $d = 3$), while Bj\"orklund~\cite{bjorklund2012,bjorklund2016} established $O^*(2^{n/2})$ and below-all-subsets bounds for perfect matchings (with fast subset convolution~\cite{bhkk2007} and Hafnian generalizations~\cite{bgq2019}). On the empirical side, zero-skipping Gray-code solvers such as SpaRyser and SkipPer~\cite{kaya2019} skip zero-product iterations without guaranteeing a tight base on regular instances. Because the pathwidth of an $m$-vertex cubic graph satisfies $pw(G) \le m/6 + o(m)$~\cite{fominhoie2006}, a $2n$-vertex 3-regular bipartite graph has pathwidth bounded near $n/3$ in the worst case, placing straightforward frontier DP near $2^{n/3} \approx 1.26^n$. Combining degree-2 contractions with min-frontier ordering pushes empirical scaling on random 3-regular instances down to $1.238^n$ (and $1.188^n$ with the repository's Or-Opt SA extension).

While approximate algorithms (such as the Jerrum--Sinclair--Vigoda Markov Chain Monte Carlo polynomial-time approximation scheme~\cite{jsv2004}) exist for non-negative matrices, \emph{exact} algebraic permanent evaluation over arbitrary finite fields $\F_p$ or complex fields $\mathbb{C}$ remains computationally demanding. Standard dense solvers fail to exploit sparsity effectively because hypercube polarization treats all coordinates symmetrically.

\paragraph{Our Contributions.}
In this paper, we present an exact, deterministic solver that achieves empirical sub-$2^n$ scaling on sparse and bounded-degree matrices, without claiming a new worst-case asymptotic theorem beyond the pathwidth and reduction framework. First, we prove that any degree-2 row or column in a bipartite graph can be contracted in $O(\deg(c_1)+\deg(c_2))$ operations ($O(1)$ when neighbor degrees are $O(1)$, and up to $O(n)$ as degrees grow under iterated contraction; algebraic cancellation to $0$ can also reduce degrees further) by algebraically merging its neighbors, reducing the matrix dimension from $n$ to $n-1$ while preserving $\Perm(A)$ identically. Second, we formulate a greedy frontier-minimizing row-ordering heuristic targeting small frontier width $fw(\pi) \approx pw(G)$ (without a worst-case guarantee of achieving $pw(G)$), which reduces the DP state space from $2^n$ to $2^{fw(\pi)}$. Finally, we validate our implementation across benchmark suites against trusted modular Glynn baselines: on random 3-regular weighted graphs, we observe an empirical scaling of $1.238^n$, solving $n = 28$ instances in $3.60$\,ms with $5{,}118$ operations (a $52{,}446\times$ reduction compared to $2^n = 268{,}435{,}456$).

\section{Iterated Exact Graph Reductions}

Let $G = (U, V, E, W)$ denote the edge-weighted bipartite graph associated with matrix $A \in \F^{n \times n}$, where $U = \{u_1, \dots, u_n\}$ indexes rows, $V = \{v_1, \dots, v_n\}$ indexes columns, and edge $(u_i, v_j)$ has weight $A_{ij} \in \F$.

\subsection{Forced Edge Peeling (Degree-1 Elimination)}

\begin{lemma}[Degree-1 Peeling]
If any row or column of $A$ is identically zero, $\Perm(A) = 0$. If row $u_r$ has degree 1 with single non-zero entry $A_{r, c} = w \neq 0$, then:
\begin{equation}
\Perm(A) = w \cdot \Perm(A_{\setminus r, \setminus c}),
\end{equation}
where $A_{\setminus r, \setminus c}$ is the $(n-1) \times (n-1)$ submatrix formed by removing row $r$ and column $c$ (with the convention $\Perm(\text{empty } 0\times 0) = 1$ for $n = 1$).
\end{lemma}

The same property holds symmetrically for any column of degree 1.

\subsection{Exact Degree-2 Series Contraction}

The central algebraic reduction is the contraction of degree-2 vertices, which eliminates paths of length 2 without branching:

\begin{theorem}[Degree-2 Series Contraction]
\label{thm:deg2_contraction}
Let row $r$ have exactly two non-zero entries: $A_{r, c_1} = w_1$ and $A_{r, c_2} = w_2$. Construct an $(n-1) \times (n-1)$ matrix $A'$ by deleting row $r$ and column $c_2$, and updating the entries of column $c_1$ for all remaining rows $i \neq r$ according to:
\begin{equation}
A'_{i, c_1} = w_1 A_{i, c_2} + w_2 A_{i, c_1}.
\end{equation}
Then:
\begin{equation}
\Perm(A) = \Perm(A').
\end{equation}
\end{theorem}

\begin{proof}
By cofactor expansion along row $r$:
\begin{equation}
\Perm(A) = w_1 \Perm(A_{\setminus r, \setminus c_1}) + w_2 \Perm(A_{\setminus r, \setminus c_2}).
\end{equation}
In both $(n-1) \times (n-1)$ matrices $A_{\setminus r, \setminus c_1}$ and $A_{\setminus r, \setminus c_2}$, all columns other than $c_1$ and $c_2$ are identical. Using invariance of $\Perm$ under arbitrary column permutations (including the cyclic shift aligning the surviving column $c_2$ in $A_{\setminus r, \setminus c_1}$ with $c_1$ in $A_{\setminus r, \setminus c_2}$ without sign change), multilinearity of the permanent with respect to its columns yields:
\begin{equation}
w_1 \Perm(A_{\setminus r, \setminus c_1}) + w_2 \Perm(A_{\setminus r, \setminus c_2}) = \Perm(\tilde{A}),
\end{equation}
where $\tilde{A}$ replaces column $c_1$ with $w_2 \cdot \mathrm{col}(c_1) + w_1 \cdot \mathrm{col}(c_2)$ and deletes column $c_2$ and row $r$. This matches the definition of $A'$ (noting that entries $w_1 A_{i,c_2} + w_2 A_{i,c_1}$ may cancel to $0$ in $\F$, further reducing degree), completing the proof.
\end{proof}

\begin{corollary}[Column Series Contraction]
Symmetrically, if column $c$ has degree 2 with entries $A_{r_1, c} = w_1$ and $A_{r_2, c} = w_2$, contracting column $c$ and merging rows $r_1, r_2$ via:
\begin{equation}
A'_{r_1, j} = w_1 A_{r_2, j} + w_2 A_{r_1, j}
\end{equation}
preserves $\Perm(A) = \Perm(A')$ over any field $\F$.
\end{corollary}

\subsection{Component Factorization and Hall Violation}

If the bipartite incidence graph decomposes into $k$ connected components $G = G_1 \sqcup \dots \sqcup G_k$ with row sets $U_m$ and column sets $V_m$, two cases arise. If $|U_m| \neq |V_m|$ for any component (a necessary component-balance violation implying a Hall violation on $U_m$ or $V_m$; Hall violations inside balanced components are not separately tested), no perfect matching can exist and $\Perm(A) = 0$ immediately. Otherwise, if $|U_m| = |V_m| = n_m$ for all $m$, the matrix can be permuted into block diagonal form, yielding
\begin{equation}
\Perm(A) = \prod_{m=1}^k \Perm(A[U_m, V_m]).
\end{equation}

\section{Min-Frontier Dynamic Programming}

When iterated degree-1 and degree-2 reductions reach a fixed point (where all remaining vertices have degree $\ge 3$), the reduced component is evaluated using frontier dynamic programming guided by a low-width path decomposition.

\subsection{Path Decomposition and the Frontier Cut}

Let $\pi = (\pi_1, \dots, \pi_n)$ be a permutation of the rows. At step $t \in \{1, \dots, n\}$, let $R_t = \{\pi_1, \dots, \pi_t\}$ be the processed rows. A column $c \in \{1, \dots, n\}$ is classified as:
\begin{itemize}[leftmargin=*]
    \item \emph{Inactive (unopened):} $c \notin \bigcup_{i=1}^t \mathrm{supp}(A_{\pi_i, :})$.
    \item \emph{Active (frontier):} $c \in \bigcup_{i=1}^t \mathrm{supp}(A_{\pi_i, :})$ and there exists some future row $t' > t$ such that $A_{\pi_{t'}, c} \neq 0$.
    \item \emph{Dead (closed):} $c \in \bigcup_{i=1}^t \mathrm{supp}(A_{\pi_i, :})$ and for all future rows $t' > t$, $A_{\pi_{t'}, c} = 0$.
\end{itemize}

The \emph{frontier width} at step $t$, denoted $fw(t)$, is the number of currently active columns. The maximum frontier width is:
\begin{equation}
fw(\pi) = \max_{1 \le t \le n} fw(t).
\end{equation}

\subsection{Greedy Min-Frontier Row Ordering}

Finding the ordering $\pi$ that minimizes $fw(\pi)$ corresponds to a one-sided (row-ordered) pathwidth variant on the bipartite incidence graph (distinct from unconstrained bipartite pathwidth, which allows interleaving row and column vertices, and likewise $\NP$-hard in general). We use a greedy heuristic that iteratively selects the next row $\pi_{t+1}$ minimizing the net change in active frontier size:
\begin{equation}
\mathrm{cost}(r) = \Delta \mathrm{Open}(r) - \Delta \mathrm{Close}(r),
\end{equation}
where $\Delta \mathrm{Open}(r)$ is the count of previously unopened columns appearing in row $r$, and $\Delta \mathrm{Close}(r)$ is the count of columns whose last non-zero appearance is in row $r$ (columns whose first and last appearance is row $r$ both open and close at $r$, contributing net change $0$).

\subsection{Frontier State Compression}

Standard Ryser DP stores all $2^n$ column subsets. In contrast, our frontier DP restricts bitmasks solely to the currently active frontier indices:
\begin{equation}
\mathrm{DP}_t : 2^{fw(t)} \to \F_p.
\end{equation}
When a column becomes dead at step $t$, all states where that column was not matched are discarded (as they cannot form a full matching), and the dead column's bit is marginalized out, compressing the active table size to $2^{fw(t)}$.

\begin{proposition}[Frontier DP Complexity]
The number of dynamic programming transitions required to evaluate a reduced matrix with path decomposition ordering $\pi$ is bounded by:
\begin{equation}
T(n) \le O\left(n \cdot d \cdot 2^{fw(\pi)}\right),
\end{equation}
where $d$ is the maximum row degree of the reduced matrix (which can exceed the input graph's degree after degree-2 contractions). When dead columns are projected on-the-fly into $\mathrm{DP}_t$, space remains $O(2^{fw(\pi)})$ (whereas materializing an intermediate table over $F_{t-1} \cup \mathrm{supp}(A_{\pi_t, :})$ before compression can use $2^{fw(\pi)+d}$ space).
\end{proposition}

\subsection{Adaptive Cascade-Branching at Degree 3}

When the graph reductions reach a fixed point where every remaining vertex has degree $\ge 3$, the solver evaluates the frontier cut width. In the reference implementation (\texttt{src/sparse\_permanent.py}), if the component size $k > 12$ and either the minimum row degree is $3$ with $fw > 9$ or $fw > 16$, the algorithm activates \emph{adaptive cascade-branching} on an edge $(r^*, c^*)$ chosen by selecting a row $r^*$ of minimum degree (breaking ties by minimum neighbor column degree) and its minimum-degree column $c^*$:
\begin{equation}
\Perm(A) = w^* \cdot \Perm(A_{\setminus r^*, \setminus c^*}) + \Perm(A_{- (r^*, c^*)}).
\end{equation}

\begin{proposition}[First-Step Cascade-Branching Recurrence]
Let $G$ be a 3-regular bipartite graph. Selecting an edge $(r^*, c^*)$ yields two subproblems:
\begin{enumerate}[leftmargin=*]
    \item \emph{Inclusion branch ($A_{\setminus r^*, \setminus c^*}$):} Removing row $r^*$ and column $c^*$ deletes 4 adjacent edges, immediately reducing 4 neighboring vertices to degree 2 (or fewer if they share edges).
    \item \emph{Exclusion branch ($A_{- (r^*, c^*)}$):} Deleting edge $(r^*, c^*)$ immediately reduces vertex $r^*$ and vertex $c^*$ to degree 2.
\end{enumerate}
In both branches, the appearance of degree-2 vertices immediately triggers Theorem~\ref{thm:deg2_contraction} series contractions. In the exclusion branch, contracting $r^*$ and $c^*$ reduces the dimension to $n - 2$ (or smaller). In the inclusion branch, when $(r^*, c^*)$ lies in no $C_4$ (e.g., girth $\ge 6$ or $C_4$-free neighborhoods), the four degree-2 neighbors contract independently to yield dimension $n - 5$ (or smaller if parallel edges or degree-1 vertices form); however, if $(r^*, c^*)$ participates in one or more $C_4$s, contracting one degree-2 neighbor merges another into a degree-3 vertex, so the inclusion branch reduces only to $n - 4$ (one $C_4$), $n - 3$ (two $C_4$s sharing $(r^*, c^*)$), or $n - 2$ (e.g., $K_{3,3}$). Under $C_4$-free neighborhoods, the local first-step recurrence $T(n) \le T(n - 5) + T(n - 2)$ has characteristic root $\lambda \approx 1.236$; moreover, contracting degree-2 vertices can merge degree-3 neighbors into vertices of degree $4$ (or higher), so deeper branching steps need not remain 3-regular (consistent with the empirical $1.303^n$ deep-branching scaling in Section~4.2).
\end{proposition}

\section{Experimental Results}

\subsection{Empirical Scaling on $d$-Regular Bipartite Graphs}

To evaluate the sub-$2^n$ scaling of the solver, we benchmarked random weighted $d$-regular bipartite graphs (constructed as random $d$-regular bipartite multigraphs/simple graphs with uniform nonzero weights in $\F_p$, $p = 2^{61}-1$; ``Generation'' labels throughout denote successive solver revisions in the accompanying repository). Averages in Tables~\ref{tab:sparse_scaling}--\ref{tab:cpp_sparse_bench} are taken over multiple random instances per size (hence fractional mean operation and branch counts; individual graph instances vary depending on expansion and cut width). In these instances, graph reductions simplify peripheral paths, adaptive cascade-branching triggers cascading degree-2 contractions, and min-frontier ordering finds narrow path decompositions.

\begin{table}[!htbp]
\centering
\caption{Empirical complexity and wall-clock execution time of \texttt{sparse\_permanent.py} (Generation 11) vs.\ dense Glynn ($2^n$) across regular bipartite graphs.}
\label{tab:sparse_scaling}
\resizebox{\linewidth}{!}{%
\begin{tabular}{lrrrrrr}
\toprule
Graph Family & $n$ & Avg DP Ops & Dense $2^n$ Ops & Ratio (Speedup) & Max Frontier & Runtime \\
\midrule
3-regular weighted & 16 & 367    & 65,536      & $178.5\times$    & 9  & 0.36 ms \\
3-regular weighted & 20 & 742    & 1,048,576   & $1,413.2\times$  & 10 & 0.62 ms \\
3-regular weighted & 24 & 1,273  & 16,777,216  & $13,183.1\times$ & 10 & 0.97 ms \\
3-regular weighted & 28 & 5,118  & 268,435,456 & $52,446.2\times$ & 13 & 3.60 ms \\
\textbf{3-regular weighted} & \textbf{36} & \textbf{17,212} & \textbf{68,719,476,736} & \textbf{3,992,532.9$\times$} & \textbf{15} & \textbf{34.3 ms} \\
\midrule
4-regular weighted & 16 & 2,596  & 65,536      & $25.2\times$     & 12 & 1.55 ms \\
4-regular weighted & 20 & 9,841  & 1,048,576   & $106.5\times$    & 14 & 6.20 ms \\
\textbf{4-regular weighted} & \textbf{24} & \textbf{32,477} & \textbf{16,777,216} & \textbf{516.6$\times$} & \textbf{16} & \textbf{19.71 ms} \\
\bottomrule
\end{tabular}%
}
\end{table}

\subsection{Effective Exponential Base Fitting}

Fitting the operation count $T(n) = C \cdot \alpha^n$ yields the following empirical bases:
\begin{itemize}[leftmargin=*]
    \item \emph{3-regular bipartite graphs:} $\alpha \approx 1.216$--$1.238$ in the frontier-DP regime ($1.2062^n$ on 0--1 unweighted instances, solving $n = 40$ in $7{,}106$ ops / $127$\,ms) and $1.303$ across deep cascade-branching up to $n = 48$. With the repository's implementation extensions (128-bit mixed-vertex DP, 0-op Boolean reachability, and Dulmage--Mendelsohn + Or-Opt SA ordering), 3-regular scaling improves further to $1.188^n$ at $n = 60$.
    \item \emph{4-regular bipartite graphs:} $\alpha \approx 1.416$ in standard frontier-DP, dropping by $46\%$ in operations to $1.314^n$ under the 128-bit mixed-vertex Dulmage--Mendelsohn engine.
    \item \emph{Sparse random and banded matrices:} 15--25\% density random matrices scale at $1.8501^n$ (vs.\ a finite-$n$ exponential fit of $2.3235^n$ to Ryser's $\Theta(n \cdot 2^n)$ baseline), while banded matrices ($w \in \{4,5,6\}$) achieve $O(n)$ linear scaling, evaluating $n = 32$ in under $1$\,ms.
\end{itemize}

\subsection{Native C++17 \& OpenMP Execution up to $n = 72$ (Generations 12--15)}

To eliminate Python interpreter overhead and evaluate exact modular arithmetic natively, we implemented the full reduction, cascade-branching, and min-frontier DP pipeline in C++17 (\texttt{src/sparse\_\-permanent.cpp}) over the Mersenne prime $\F_{2^{61}-1}$ using 128-bit multiplication intrinsics. In Generation~15, raising the flat-array zero-hash DP cutoff to \texttt{fw} $\le 15$ eliminates over $95\%$ of recursive branches, speeding up $n = 56$ by $43.7\times$ and $n = 64$ by $16.3\times$, and solving exact $72 \times 72$ weighted permanents in $2.12$\,seconds (2 threads) and $580.49$\,ms (8 OpenMP threads). Table~\ref{tab:cpp_sparse_bench} reports wall-clock runtimes for Generation~15.

\begin{table}[!htbp]
\centering
\caption{Generation 15 Native C++17 + OpenMP (\texttt{src/sparse\_permanent.cpp}) exact 3-regular weighted permanent benchmarks ($n = 48\text{--}72$) modulo $2^{61}-1$.}
\label{tab:cpp_sparse_bench}
\begin{tabular*}{\linewidth}{@{\extracolsep{\fill}}rrrrrr}
\toprule
$n$ & Avg Ops & Dense $2^n$ Ops & Ratio (Speedup) & Branches & Gen 15 Time \\
\midrule
48 & 228,830     & $2.81 \times 10^{14}$ & $1.23 \times 10^9\times$     & 5.5     & 4.90 ms \\
56 & 529,837     & $7.21 \times 10^{16}$ & $1.36 \times 10^{11}\times$  & 20.0    & 9.83 ms \\
64 & 34,599,126  & $1.84 \times 10^{19}$ & $5.33 \times 10^{11}\times$  & 1,550.0 & 599.13 ms \\
68 & 143,534,125 & $2.95 \times 10^{20}$ & $2.06 \times 10^{12}\times$  & 7,962.0 & 2.76 s \\
\textbf{72} & \textbf{116,509,326} & \textbf{4.72 $\times$ 10$^{21}$} & \textbf{4.05 $\times$ 10$^{13}\times$} & \textbf{6,057.0} & \textbf{2.12 s} \\
\bottomrule
\end{tabular*}
\end{table}

At $n = 72$, where dense $2^{72}$ requires $4.72 \times 10^{21}$ ($4.72$ sextillion) operations, the Generation~15 OpenMP + flat-array C++17 engine computes the exact permanent in $2.12$\,s (2 threads) or $580.49$\,ms (8 threads; $n=56$ in $3.93$\,ms, $n=64$ in $192.53$\,ms, $n=68$ in $825.31$\,ms) with $116{,}509{,}326$ operations ($4.05 \times 10^{13}\times$ fewer operations), verified permutation-invariant across all sizes.

\subsection{GPU Acceleration on NVIDIA Tesla T4 (CUDA)}

To evaluate GPU-batched frontier DP, we ported the engine to CUDA (\texttt{colab/sparse\_\-permanent\_\-cuda.cu}) with chunked VRAM batching (\texttt{MAX\_BATCH} $= 2048$, $\approx 4$\,GB cap) and 64-bit Mersenne modular intrinsics. On an NVIDIA Tesla T4 (SM 7.5, 15.6 GB), it passed 5/5 Glynn self-tests ($n=6\text{--}14$) and solved 3-regular weighted matrices through $n = 72$:
\begin{itemize}[leftmargin=*]
    \item $n = 32$: $0.88$\,ms ($0$ branches); $n = 40$: $1.34$\,ms ($0$ branches);
    \item $n = 48$: $5.07$\,ms ($35$ branches); $n = 56$: $74.23$\,ms ($1{,}185$ branches);
    \item $n = 64$: $1.40$\,s ($16{,}452$ branches); $n = 72$: $20.14$\,s ($196{,}343$ branches).
\end{itemize}
Adding $32$\,KB on-chip \texttt{\_\_shared\_\_} DP for \texttt{fw} $\le 11$ (Gen~17) reduced $n=72$ to $19.04$\,s; adding \texttt{\_\_launch\_bounds\_\_(256)} and atomic-free bijective dead-column compression (Gen~18) gave $72.36$\,ms ($n=56$), $1.37$\,s ($n=64$), and $19.80$\,s ($n=72$).

\takeaway{Sub-$2^n$ Scaling on Sparse Graphs}{While Part~I established $O(2^n)$ non-scalar multiplications for linear polarization on dense matrices, structured and bounded-degree matrices exhibit much smaller effective cut widths. By interleaving degree-2 series contractions with frontier-compressed dynamic programming, we observe empirical sub-$2^n$ scaling ($1.238^n$ for random $d=3$ instances) without sacrificing exactness, complementing proved worst-case bounds~\cite{servediowan2005,fominhoie2006,izumiwadayama2012,bjorklund2012,bjorklund2016}.}

\section{Conclusion and Future Outlook}

We have presented an exact, verified permanent solver that achieves empirical exponential bases strictly below $2$ on bounded-degree and sparse bipartite benchmarks over arbitrary fields. By coupling exact degree-2 series contractions with min-frontier path decomposition, cascade-branching, and the repository's Generation~9 Butterfly OptTree dense-fallback extension, the algorithm achieves up to a $4.05 \times 10^{13}\times$ reduction in operations at $n = 72$.

In Generations~16--20, the repository was extended along five additional engineering directions beyond the core bipartite algorithm analyzed in Sections~2--3. First, for complex-valued Boson Sampling (\texttt{src/boson\_\-sampling\_\-complex.cpp}), compensated \texttt{long double} accumulation solves sparse 3-regular $n = 60$ unitary permanents in $192.15$\,ms (though general dense $n=60$ unitaries still require $\sim 2^{60}$ operations, and finite-precision compensated \texttt{long double} does not guarantee machine-precision relative accuracy under catastrophic cancellation in general complex or unitary matrices). Second, for non-bipartite Gaussian Boson Sampling graphs, an exact sub-$2^n$ Sparse Hafnian solver (\texttt{src/sparse\_\-hafnian.cpp}, cf.~\cite{bgq2019}) with cascade-branching and OpenMP solves 3-regular $n = 72$ instances in $5.09$\,ms and scales to $n = 96$ in $112.64$\,ms (as well as 4-regular $n = 32$ in $0.551$\,ms). Third, the batched CUDA GPU kernel (\texttt{colab/sparse\_\-permanent\_\-cuda.cu}) with $32$\,KB on-chip shared-memory DP, \texttt{\_\_launch\_\-bounds\_\_(256)}, and atomic-free compression solves $n = 64$ in $1.37$\,s and $n = 72$ in $19.04$\,s--$19.80$\,s on a Tesla T4. Fourth, combining Or-Opt, column-closing, and remaining-degree frontier ordering with \texttt{fw} $\le 16$ flat-array DP in \texttt{src/sparse\_\-permanent.cpp} scales 3-regular weighted permanents to $n = 80$ ($744.24$\,M ops, $1.62 \times 10^{15}\times$ fewer than $2^{80}$, $21.77$\,s) and 4-regular (\texttt{--4reg}) instances to $n = 28$ in $94{,}364$ ops (\texttt{br} $= 0$) and $n = 36$ in $42.48$\,ms ($1.80$\,M ops). Finally, in Lean~4 with Mathlib4 (\texttt{lean/Permanent\-Proofs.lean}), full Laplace cofactor expansion and row-collapse theorems (\texttt{permanent\_\-succ\_\-row\_\-zero}, \texttt{permanent\_\-deg1\_\-peeling}, and \texttt{permanent\_\-deg2\_\-collapse}) typecheck cleanly with zero \texttt{sorry}s.

\section*{Reproducibility Package}
All source code, C++17/CUDA kernels, Lean~4 formalizations, and verification suites are publicly accessible and reproducible at: \\
\url{https://github.com/3DBarath/matrix_permanent_circuit_discovery}.

\section*{AI Assistance Disclosure}
AI coding and writing assistants were used during algorithm discovery, C++17/CUDA/Lean~4 implementation, and manuscript drafting. All mathematical proofs, bit-exact Glynn/brute-force checks, hardware benchmarks, and Lean~4 + Mathlib4 formalizations were verified deterministically by the author.

\begin{thebibliography}{99}
\bibitem{valiant1979} L.~G.~Valiant, ``The complexity of computing the permanent,'' \emph{Theoretical Computer Science}, vol.~8, no.~2, pp.~189--201, 1979.
\bibitem{ryser1963} H.~J.~Ryser, \emph{Combinatorial Mathematics}, Carus Mathematical Monographs, No.~14, Mathematical Association of America, 1963.
\bibitem{glynn2010} D.~G.~Glynn, ``The permanent of a square matrix,'' \emph{European Journal of Combinatorics}, vol.~31, no.~7, pp.~1887--1891, 2010.
\bibitem{part1_paper} Barath~B, ``Compressing Glynn's Formula: Exact $O(2^n)$-Multiplication Circuits for the Matrix Permanent via Degree-Preserving Dual-Ball Annihilation and M\"obius--Zeta Butterflies,'' \emph{Preprint}, 2026.
\bibitem{aaronson2011} S.~Aaronson and A.~Arkhipov, ``The computational complexity of linear optics,'' in \emph{Proceedings of the 43rd Annual ACM Symposium on Theory of Computing (STOC)}, pp.~333--342, 2011.
\bibitem{jsv2004} M.~Jerrum, A.~Sinclair, and E.~Vigoda, ``A polynomial-time approximation algorithm for the permanent of a matrix with non-negative entries,'' \emph{Journal of the ACM}, vol.~51, no.~4, pp.~671--697, 2004.
\bibitem{servediowan2005} R.~A.~Servedio and A.~Wan, ``Computing sparse permanents faster,'' \emph{Information Processing Letters}, vol.~96, no.~3, pp.~89--92, 2005.
\bibitem{fominhoie2006} F.~V.~Fomin and K.~H\o ie, ``Pathwidth of cubic graphs and exact algorithms,'' \emph{Information Processing Letters}, vol.~97, no.~5, pp.~191--196, 2006.
\bibitem{bhkk2007} A.~Bj\"orklund, T.~Husfeldt, P.~Kaski, and M.~Koivisto, ``Fourier meets M\"obius: fast subset convolution,'' in \emph{Proceedings of the 39th Annual ACM Symposium on Theory of Computing (STOC)}, pp.~67--74, 2007.
\bibitem{izumiwadayama2012} T.~Izumi and T.~Wadayama, ``A new direction for counting perfect matchings,'' in \emph{Proceedings of the 53rd Annual IEEE Symposium on Foundations of Computer Science (FOCS)}, pp.~591--598, 2012.
\bibitem{bjorklund2012} A.~Bj\"orklund, ``Counting perfect matchings as fast as Ryser,'' in \emph{Proceedings of the 23rd Annual ACM-SIAM Symposium on Discrete Algorithms (SODA)}, pp.~914--921, 2012.
\bibitem{bjorklund2016} A.~Bj\"orklund, ``Below all subsets for some permutational counting problems,'' in \emph{15th Scandinavian Symposium and Workshops on Algorithm Theory (SWAT)}, LIPIcs, vol.~53, pp.~17:1--17:11, 2016.
\bibitem{bgq2019} A.~Bj\"orklund, B.~Gupt, and N.~Quesada, ``A faster hafnian formula for complex matrices and its benchmarking on a supercomputer,'' \emph{ACM Journal of Experimental Algorithmics}, vol.~24, art.~1.11, pp.~1--17, 2019.
\bibitem{kaya2019} K.~Kaya, ``Parallel algorithms for computing sparse matrix permanents,'' \emph{Turkish Journal of Electrical Engineering and Computer Sciences}, vol.~27, no.~6, pp.~4284--4297, 2019.
\end{thebibliography}

\end{document}
